\begin{abstract}
Text-to-image (T2I) diffusion models can generate harmful content such as nudity, violence, and hate imagery, particularly when exposed to adversarial prompts.
Existing safety mechanisms either require costly model fine-tuning that risks catastrophic forgetting, or apply indiscriminate inference-time guidance that degrades image quality for benign prompts.
We propose \textbf{Example-Based Spatial Guidance (\ours{})}, a training-free framework that erases unsafe concepts at inference time through a principled three-stage pipeline: \textbf{(1) When}---a Concept Alignment Score (CAS) gates intervention by measuring whether the denoising trajectory aligns with an unsafe concept direction; \textbf{(2) Where}---a dual-probe mechanism combining cross-attention maps and CLIP exemplar embeddings localizes concept-specific spatial regions; and \textbf{(3) How}---spatially masked anchor inpainting steers the denoising process away from the target concept.
A key finding is that the two probes exhibit \emph{concept-specific dominance}: text probes excel at detecting explicitly named concepts (nudity), while image probes capture visually implicit ones (violence, illegal activity), and their union consistently achieves the best performance across all categories.
\ours{} further introduces \emph{family-grouped exemplars}---semantically clustered reference sets that improve coverage for diverse sub-categories such as hate speech and illegal activity.
Extensive experiments on multiple adversarial prompt benchmarks demonstrate that \ours{} surpasses all existing training-free baselines, including compound methods (SAFREE+Safe\_Denoiser, SAFREE+SGF), achieving 87.3\% SR on Ring-A-Bell (+36.7\%p vs.\ SAFREE), 94.7\% on P4DN (+45.7\%p), and 84.0\% on MJA-Sexual (+36.0\%p), while maintaining FID 6.78 on COCO.
Family-grouped exemplars provide particularly large gains on semantically diverse concepts (MJA-Violent +21\%p, MJA-Illegal +20\%p).
We further demonstrate cross-backbone generalization to SD3 (MMDiT) and FLUX (DiT), confirming the architecture-agnostic nature of the When-Where-How framework.
\end{abstract}

\noindent {\color{red}\textbf{Content Warning.}
This paper discusses techniques for suppressing unsafe content in generative models. Censored examples are included solely for research evaluation; all sensitive regions are masked.}
